\documentclass[12pt]{report}
\def\TITULO{Informe Técnico Final}
\def\subtitulo{Calendario de Recordatorios}
\def\autora{Franco Garrido}
\def\correoa{francogamer950@gmail.com}

\def\asignatura{Taller de Ingeniería Informática}
\def\campus{Campus Osorno}
\def\carrera{Ingeniería Civil en Informática}
\input{Macros}

\lstset{language=PHP}

\begin{document}

\maketitle

\cleardoublepage
\pagenumbering{roman}
\setcounter{page}{1}

\tableofcontents
\listoffigures

\cleardoublepage
\pagenumbering{arabic}
\setcounter{page}{1}

\chapter{Introducción}

El informe técnico tiene como objetivo documentar el desarrollo y funcionamiento de la aplicación web \textbf{Calendario de Recordatorios}, un sistema enfocado en el manejo de recordatorios mediante acciones CRUD (crear, consultar, modificar y eliminar). La aplicación permite al usuario crear, consultar, modificar y eliminar registros de una forma ordenada a partir de una interfaz web.

El proyecto alberga un sistema de autenticación que permite el inicio de sesión mediante usuario y contraseña así como registrar nuevos usuarios. Los datos están almacenados en una base de datos que asocia los recordatorios con el usuario correspondiente.

Asimismo, la aplicación tiene un menú principal que distingue entre las operaciones principales de la aplicación: crear, consultar, modificar, eliminar y cerrar sesión. Esta estructura permite el acceso directo de cada operación solicitada por el usuario.

Un elemento destacado que se ha incorporado al proyecto es la implementación de una bitácora de eventos, encargada de registrar las acciones que son llevadas a cabo por el propio usuario dentro de la aplicación. Esta bitácora permite guardar la siguiente información: fecha y hora, usuario, tipo de evento, detalle de la operacion y la IP o el host del cliente.

\chapter{Desarrollo}

\section{Deploy de la Aplicación}

\begin{figure}[H]
    \centering
    \includegraphics[width=1\linewidth]{images/screenshots/deploy_aplicacion.png}
    \caption{Interfaz Menú Principal de la Aplicación}
    \label{fig:docker-up}
\end{figure}

Menú principal de la aplicación, aqui se pueden visualizar todas las tareas registradas en todos los estados, un filtro por estado, las tareas pendientes cercanas, unas tarjetas con la cantidad total de de tareas además de la cantidad de tareas por estados y por ultimo un menú para acceder a operaciones CRUD, junto al cierre de sesión.

\section{Interfaces de autenticación}

\subsection{Inicio de sesión}

\begin{figure}[H]
    \centering
    \includegraphics[width=0.7\linewidth]{images/screenshots/login.png}
    \caption{Interfaz de inicio de sesión}
    \label{fig:login}
\end{figure}

La interfaz de inicio de sesión permite que un usuario registrado acceda al sistema mediante nombre de usuario y contraseña. La contraseña se valida desde PHP utilizando el hash almacenado en la base de datos.

\subsection{Registro de usuario}

\begin{figure}[H]
    \centering
    \includegraphics[width=0.7\linewidth]{images/screenshots/registro_usuario.png}
    \caption{Interfaz de registro de usuario}
    \label{fig:registro}
\end{figure}

La interfaz de registro permite crear un nuevo usuario dentro del sistema. Al guardar el usuario, la contraseña no se almacena en texto plano, sino que se guarda en la columna \texttt{contrasena} utilizando un hash generado por PHP.

\section{Menú principal de operaciones CRUD}

\begin{figure}[H]
    \centering
    \includegraphics[width=0.3\linewidth]{images/screenshots/menu_principal.png}
    \caption{Menú principal con opciones CRUD}
    \label{fig:menu-principal}
\end{figure}

El menú principal concentra las opciones: \textbf{Crear}, \textbf{Consultar}, \textbf{Modificar}, \textbf{Eliminar} y \textbf{Cerrar sesión}. Desde este menú el usuario puede acceder a cada operación de manera separada.

\section{Operaciones CRUD}

\subsection{Crear registro}

\begin{figure}[H]
    \centering
    \includegraphics[width=0.9\linewidth]{images/screenshots/crear_recordatorio.png}
    \caption{Formulario para crear un recordatorio}
    \label{fig:crear}
\end{figure}

Esta interfaz permite registrar un nuevo recordatorio indicando sus datos principales, como título, descripción, fecha y hora. Al guardar el formulario, el registro queda asociado al usuario que inició sesión.

\subsection{Consultar registros}

\begin{figure}[H]
    \centering
    \includegraphics[width=0.9\linewidth]{images/screenshots/consultar_recordatorios.png}
    \caption{Consulta y listado de recordatorios}
    \label{fig:consultar}
\end{figure}

La vista de consulta muestra los recordatorios registrados por el usuario. También permite revisar el estado de cada tarea y aplicar filtros para visualizar registros específicos.

\subsection{Modificar registro}

\begin{figure}[H]
    \centering
    \includegraphics[width=0.9\linewidth]{images/screenshots/modificar_recordatorio.png}
    \caption{Selección y modificación de un recordatorio}
    \label{fig:modificar}
\end{figure}

La operación de modificación permite seleccionar un recordatorio existente y actualizar sus datos.

\subsection{Eliminar registro}

\begin{figure}[H]
    \centering
    \includegraphics[width=0.9\linewidth]{images/screenshots/eliminar_recordatorio.png}
    \caption{Eliminación de un recordatorio}
    \label{fig:eliminar}
\end{figure}

La interfaz de eliminación permite borrar un recordatorio específico. Esta acción se registra en la bitácora para mantener trazabilidad sobre los cambios realizados dentro del sistema.

\section{Bitácora}

\begin{figure}[H]
    \centering
    \includegraphics[width=0.9\linewidth]{images/screenshots/base_datos_log.png}
    \caption{Modelo de datos de usuarios, recordatorios y bitácora}
    \label{fig:modelo-datos}
\end{figure}

La bitácora registra los eventos principales de la aplicación. Cada evento considera la fecha y hora, el nombre de usuario, el tipo de acción realizada, el detalle de la operación y la IP o host del cliente. Esto permite mantener un historial de uso del sistema y facilita la revisión de acciones realizadas por los usuarios.

\chapter{Conclusión}

El desarrollo de la aplicación \textbf{Calendario de Recordatorios} permitió implementar un sistema web funcional para la administración de recordatorios mediante operaciones CRUD. La aplicación cumple con las funciones principales de crear, consultar, modificar y eliminar registros, entregando al usuario una forma ordenada de gestionar sus tareas.

La incorporación del inicio de sesión y el registro de usuarios permite controlar el acceso al sistema, asegurando que cada usuario pueda trabajar con sus propios datos. Además, el uso de una base de datos permite almacenar de manera persistente tanto los usuarios como los recordatorios registrados.

El menú principal facilita la navegación dentro de la aplicación, ya que separa cada operación en una opción específica. Esto permite que el usuario acceda directamente a la acción que desea realizar, manteniendo una estructura clara y fácil de utilizar.

Finalmente, la bitácora de eventos permite registrar las acciones relevantes realizadas dentro del sistema, entregando trazabilidad sobre el uso de la aplicación.

\end{document}
